\documentclass[twocolumn,a4j,dvipdfmx]{jsarticle}
\usepackage[top=20mm,bottom=30mm,left=20mm,right=20mm]{geometry}
\usepackage[dvipdfmx]{graphicx}
\usepackage{listings,jlisting}
\setlength{\columnsep}{2zw}
\begin{document}

\title{\LaTeX 課題---回路の過渡応答}
\author{2022311042 \quad 神野友哉}
\date{2022年6月18日}
\maketitle

\section{\Large LCR直列回路の方程式}
図\ref{fig:circuit}のように電源$V_i$，抵抗\textit{R}，コイル\textit{L}，コンデン\quad サ\textit{C}を直列に接続した回路を考える．
\vskip\baselineskip
\begin{figure}[htb]
  \centering
  \includegraphics[width=1\linewidth]{/home/shinno/circuit.eps}
  \caption{回路}
  \label{fig:circuit}
\end{figure}

コンデンサの電荷を\textit{Q}とすれば，\textit{Q}に関する次の微分方程式が得られる\cite{bibitteru}．
\begin{equation}
  L\frac{d^2Q}{dt^2}+R\frac{dQ}{dt}+\frac{1}{C}Q=V\label{eq:1}
\end{equation}
\quad\textit{V}が直流電源(または交流電源)である場合には，十分に時間が立つと\textit{Q}は一定値(または交流電源と同じ周波数の正弦関数)になる．そのような状態を定常状態といい，定常状態に至るまでの状態を過渡状態という．
\section{\Large 回路の微分方程式の解}
図\ref{fig:circuit}の回路で電源は直流であるものとする．微分方程式(\ref{eq:1}) の解を求めるために，$\lambda$に関する2次方程式
  \begin{equation}
      L\lambda^2+R\lambda+\frac{1}{C}=0\label{eq:2}
  \end{equation}
を考える．(\ref{eq:2}) の解の種類によって微分方程式(\ref{eq:1}) の解は次のようになる．ただし，\textit{A}，\textit{B}は初期条件によって決まる定数である．
\vskip\baselineskip
\begin{itemize}
\setlength{\leftskip}{1zw}  
\item (\ref{eq:2})が２つの実数解$\lambda=-\gamma_1,-\gamma_2$を持つとき
  \begin{equation}
    Q(t)=CV+Ae^{-\gamma_1t}+Be^{-\gamma_2t}\label{eq:3}
  \end{equation}
\item (\ref{eq:2})が二重解$\lambda=-\gamma$を持つとき
  \begin{equation}
    Q(t)=CV+Ae^{-\gamma t}+Bte^{-\gamma t}\label{eq:4}
  \end{equation}
\item(\ref{eq:2})が２つの複素数解$\lambda=-\gamma \pm j\omega$を持つとき\footnote{工学分野ではしばしば虚数単位を\textit{i}の代わりに\textit{j}で表す．}
  \begin{equation}
    Q(t)=CV+e^{-\gamma t}(A\mathrm{sin}{(\omega t)}+B\mathrm{cos}{(\omega t)})\label{eq:5}
  \end{equation}
\end{itemize}
\section{\Large 例題}
図\ref{fig:circuit}の回路において，$R=35.4\mathrm{\,\Omega}，L=230\mathrm{\,mH}，C=14.0\,\mu \mathrm{F}，V=1.5\mathrm{\,V}$とする．このとき，２次方程式(\ref{eq:2})の解は，$\lambda=-77.0\pm 551.9j$である．式(\ref{eq:5})において，時刻$t=0$における初期条件を$Q=0，\frac{dQ}{dt}=0$とすれば，コンデンサの両端の電圧$V_C=\frac{Q}{C}$は
\[
  V_C=1.5(1+e^{-77.0t}(-0.139\mathrm{sin}({551.9t})-\mathrm{cos}({551.9t})))
\]
となる．これを図\ref{fig:Vol}に示す．十分な時間が立つと$V_C$は一定値1.5Vに近づくことが分かる．
\vskip1\baselineskip
\begin{figure}[htb]
  \centering
  \includegraphics[width=0.8\linewidth]{/home/shinno/Voltage.eps}
  \caption{コンデンサ両端の電圧}
  \label{fig:Vol}
\end{figure}
\begin{thebibliography}{1}
  \bibitem{bibitteru} V.D.バージャー，M.G.オルソン著，戸田，田上訳，力学，培風館，東京，1975.
\end{thebibliography}
\newpage
\title
\onecolumn
  \lstinputlisting[language={Tex}]{cplitetask2.tex}
\end{document}

